
        You are a research assistant AI tasked with generating a scientific paper based on provided literature. Follow these steps:
    1. Analyze the given References.
    2. Identify gaps in existing research to establish the motivation for a new study.
    3. Propose a main idea for a new research work.
    4. Write the paper's main content in LaTeX format, including all the following parts, please must have tables:
    - Title
    - Abstract
    - Introduction
    - Related Work
    - Methods/Experiment
    - Results with tables
    - Discussion
    - Conclusion
    - Contributions
    Ensure all content is original, academically rigorous, and follows standard scientific writing conventions.Please make all decision by yourself,do not ask for any instructions

        Here are the references to be used for generating the paper.
        You MUST base the paper on these references and integrate them naturally.

        References:
        --------------------
        @misc{hossain2026llmsbenefituseritem,
      title={Do LLMs Benefit from User and Item Embeddings in Recommendation Tasks?}, 
      author={Mir Rayat Imtiaz Hossain and Leo Feng and Leonid Sigal and Mohamed Osama Ahmed},
      year={2026},
      eprint={2601.04690},
      archivePrefix={arXiv},
      primaryClass={cs.LG},
      url={https://arxiv.org/abs/2601.04690},
      abstract = {Large Language Models (LLMs) have emerged as promising recommendation systems, offering novel ways to model user preferences through generative approaches. However, many existing methods often rely solely on text semantics or incorporate collaborative signals in a limited manner, typically using only user or item embeddings. These methods struggle to handle multiple item embeddings representing user history, reverting to textual semantics and neglecting richer collaborative information. In this work, we propose a simple yet effective solution that projects user and item embeddings, learned from collaborative filtering, into the LLM token space via separate lightweight projector modules. A finetuned LLM then conditions on these projected embeddings alongside textual tokens to generate recommendations. Preliminary results show that this design effectively leverages structured user-item interaction data, improves recommendation performance over text-only LLM baselines, and offers a practical path for bridging traditional recommendation systems with modern LLMs.}
}@misc{liang2026lexicalizedconstituencyparsingmiddle,
      title={Lexicalized Constituency Parsing for Middle Dutch: Low-resource Training and Cross-Domain Generalization}, 
      author={Yiming Liang and Fang Zhao},
      year={2026},
      eprint={2601.07008},
      archivePrefix={arXiv},
      primaryClass={cs.CL},
      url={https://arxiv.org/abs/2601.07008},
      abstract = {Recent years have seen growing interest in applying neural networks and contextualized word embeddings to the parsing of historical languages. However, most advances have focused on dependency parsing, while constituency parsing for low-resource historical languages like Middle Dutch has received little attention. In this paper, we adapt a transformer-based constituency parser to Middle Dutch, a highly heterogeneous and low-resource language, and investigate methods to improve both its in-domain and cross-domain performance. We show that joint training with higher-resource auxiliary languages increases F1 scores by up to 0.73, with the greatest gains achieved from languages that are geographically and temporally closer to Middle Dutch. We further evaluate strategies for leveraging newly annotated data from additional domains, finding that fine-tuning and data combination yield comparable improvements, and our neural parser consistently outperforms the currently used PCFG-based parser for Middle Dutch. We further explore feature-separation techniques for domain adaptation and demonstrate that a minimum threshold of approximately 200 examples per domain is needed to effectively enhance cross-domain performance.}
}@misc{liu2026highrankstructuredmodulationparameterefficient,
      title={High-Rank Structured Modulation for Parameter-Efficient Fine-Tuning}, 
      author={Yongkang Liu and Xing Li and Mengjie Zhao and Shanru Zhang and Zijing Wang and Qian Li and Shi Feng and Feiliang Ren and Daling Wang and Hinrich Schütze},
      year={2026},
      eprint={2601.07507},
      archivePrefix={arXiv},
      primaryClass={cs.CL},
      url={https://arxiv.org/abs/2601.07507},
      abstract = {As the number of model parameters increases, parameter-efficient fine-tuning (PEFT) has become the go-to choice for tailoring pre-trained large language models. Low-rank Adaptation (LoRA) uses a low-rank update method to simulate full parameter fine-tuning, which is widely used to reduce resource requirements. However, decreasing the rank encounters challenges with limited representational capacity when compared to full parameter fine-tuning. We present \\textbf\{SMoA\}, a high-rank \\textbf\{S\}tructured \\textbf\{MO\}dulation \\textbf\{A\}dapter that uses fewer trainable parameters while maintaining a higher rank, thereby improving the model's representational capacity and offering improved performance potential. The core idea is to freeze the original pretrained weights and selectively amplify or suppress important features of the original weights across multiple subspaces. The subspace mechanism provides an efficient way to increase the capacity and complexity of a model. We conduct both theoretical analyses and empirical studies on various tasks. Experiment results show that SMoA outperforms LoRA and its variants on 10 tasks, with extensive ablation studies validating its effectiveness.}
}@misc{liu2026teachdiffusionlanguagemodels,
      title={Teach Diffusion Language Models to Learn from Their Own Mistakes}, 
      author={Liming Liu and Binxuan Huang and Xin Liu and Bing Yin and Tuo Zhao},
      year={2026},
      eprint={2601.06428},
      archivePrefix={arXiv},
      primaryClass={cs.LG},
      url={https://arxiv.org/abs/2601.06428},
      abstract = {Masked Diffusion Language Models (DLMs) achieve significant speed by generating multiple tokens in parallel. However, this parallel sampling approach, especially when using fewer inference steps, will introduce strong dependency errors and cause quality to deteriorate rapidly as the generation step size grows. As a result, reliable self-correction becomes essential for maintaining high-quality multi-token generation. To address this, we propose Decoupled Self-Correction (DSC), a novel two-stage methodology. DSC first fully optimizes the DLM's generative ability before freezing the model and training a specialized correction head. This decoupling preserves the model's peak SFT performance and ensures the generated errors used for correction head training are of higher quality. Additionally, we introduce Future-Context Augmentation (FCA) to maximize the correction head's accuracy. FCA generalizes the error training distribution by augmenting samples with ground-truth tokens, effectively training the head to utilize a richer, future-looking context. This mechanism is used for reliably detecting the subtle errors of the high-fidelity base model. Our DSC framework enables the model, at inference time, to jointly generate and revise tokens, thereby correcting errors introduced by multi-token generation and mitigating error accumulation across steps. Experiments on mathematical reasoning and code generation benchmarks demonstrate that our approach substantially reduces the quality degradation associated with larger generation steps, allowing DLMs to achieve both high generation speed and strong output fidelity.}
}@misc{moore2026ncbenchllmbenchmarkevaluating,
      title={NC-Bench: An LLM Benchmark for Evaluating Conversational Competence}, 
      author={Robert J. Moore and Sungeun An and Farhan Ahmed and Jay Pankaj Gala},
      year={2026},
      eprint={2601.06426},
      archivePrefix={arXiv},
      primaryClass={cs.CL},
      url={https://arxiv.org/abs/2601.06426},
      abstract = {The Natural Conversation Benchmark (NC-Bench) introduce a new approach to evaluating the general conversational competence of large language models (LLMs). Unlike prior benchmarks that focus on the content of model behavior, NC-Bench focuses on the form and structure of natural conversation. Grounded in the IBM Natural Conversation Framework (NCF), NC-Bench comprises three distinct sets. The Basic Conversation Competence set evaluates fundamental sequence management practices, such as answering inquiries, repairing responses, and closing conversational pairs. The RAG set applies the same sequence management patterns as the first set but incorporates retrieval-augmented generation (RAG). The Complex Request set extends the evaluation to complex requests involving more intricate sequence management patterns. Each benchmark tests a model's ability to produce contextually appropriate conversational actions in response to characteristic interaction patterns. Initial evaluations across 6 open-source models and 14 interaction patterns show that models perform well on basic answering tasks, struggle more with repair tasks (especially repeat), have mixed performance on closing sequences, and find complex multi-turn requests most challenging, with Qwen models excelling on the Basic set and Granite models on the RAG set and the Complex Request set. By operationalizing fundamental principles of human conversation, NC-Bench provides a lightweight, extensible, and theory-grounded framework for assessing and improving the conversational abilities of LLMs beyond topical or task-specific benchmarks.}
}@misc{chawla2026lessonsfieldadaptablelifecycle,
      title={Lessons from the Field: An Adaptable Lifecycle Approach to Applied Dialogue Summarization}, 
      author={Kushal Chawla and Chenyang Zhu and Pengshan Cai and Sangwoo Cho and Scott Novotney and Ayushman Singh and Jonah Lewis and Keasha Safewright and Alfy Samuel and Erin Babinsky and Shi-Xiong Zhang and Sambit Sahu},
      year={2026},
      eprint={2601.08682},
      archivePrefix={arXiv},
      primaryClass={cs.CL},
      url={https://arxiv.org/abs/2601.08682},
      abstract = {Summarization of multi-party dialogues is a critical capability in industry, enhancing knowledge transfer and operational effectiveness across many domains. However, automatically generating high-quality summaries is challenging, as the ideal summary must satisfy a set of complex, multi-faceted requirements. While summarization has received immense attention in research, prior work has primarily utilized static datasets and benchmarks, a condition rare in practical scenarios where requirements inevitably evolve. In this work, we present an industry case study on developing an agentic system to summarize multi-party interactions. We share practical insights spanning the full development lifecycle to guide practitioners in building reliable, adaptable summarization systems, as well as to inform future research, covering: 1) robust methods for evaluation despite evolving requirements and task subjectivity, 2) component-wise optimization enabled by the task decomposition inherent in an agentic architecture, 3) the impact of upstream data bottlenecks, and 4) the realities of vendor lock-in due to the poor transferability of LLM prompts.}
}@misc{taneja2026gpuacceleratedint8quantizationkv,
      title={GPU-Accelerated INT8 Quantization for KV Cache Compression in Large Language Models}, 
      author={Maanas Taneja and Purab Shingvi},
      year={2026},
      eprint={2601.04719},
      archivePrefix={arXiv},
      primaryClass={cs.LG},
      url={https://arxiv.org/abs/2601.04719},
      abstract = {The key-value (KV) cache in large language models presents a significant memory bottleneck during inference, growing linearly with sequence length and often exceeding the memory footprint of model weights themselves. We implement and evaluate GPU-accelerated INT8 quantization for KV cache compression, achieving 4$\\times$ memory reduction with minimal accuracy degradation. We develop four CUDA kernel variants -- naive, tiled, coarsened, and vectorized -- and benchmark them across realistic workload sizes up to 1 billion elements. Our vectorized kernel achieves up to 1,694$\\times$ speedup over CPU baselines while maintaining reconstruction error below 0.004 and attention score error below 0.1 even for 8K-dimensional heads. These results demonstrate that INT8 quantization provides a practical approach for reducing memory pressure in LLM inference with negligible computational overhead (6--58ms) and minimal impact on downstream model behavior}
}@misc{wei2026ciragconstructionintegrationretrievaladaptive,
      title={CIRAG: Construction-Integration Retrieval and Adaptive Generation for Multi-hop Question Answering}, 
      author={Zili Wei and Xiaocui Yang and Yilin Wang and Zihan Wang and Weidong Bao and Shi Feng and Daling Wang and Yifei Zhang},
      year={2026},
      eprint={2601.06799},
      archivePrefix={arXiv},
      primaryClass={cs.CL},
      url={https://arxiv.org/abs/2601.06799},
      abstract = {Triple-based Iterative Retrieval-Augmented Generation (iRAG) mitigates document-level noise for multi-hop question answering. However, existing methods still face limitations: (i) greedy single-path expansion, which propagates early errors and fails to capture parallel evidence from different reasoning branches, and (ii) granularity-demand mismatch, where a single evidence representation struggles to balance noise control with contextual sufficiency. In this paper, we propose the Construction-Integration Retrieval and Adaptive Generation model, CIRAG. It introduces an Iterative Construction-Integration module that constructs candidate triples and history-conditionally integrates them to distill core triples and generate the next-hop query. This module mitigates the greedy trap by preserving multiple plausible evidence chains. Besides, we propose an Adaptive Cascaded Multi-Granularity Generation module that progressively expands contextual evidence based on the problem requirements, from triples to supporting sentences and full passages. Moreover, we introduce Trajectory Distillation, which distills the teacher model's integration policy into a lightweight student, enabling efficient and reliable long-horizon reasoning. Extensive experiments demonstrate that CIRAG achieves superior performance compared to existing iRAG methods.}
}@misc{winata2026largelanguagemodelsunderstand,
      title={Can Large Language Models Understand, Reason About, and Generate Code-Switched Text?}, 
      author={Genta Indra Winata and David Anugraha and Patrick Amadeus Irawan and Anirban Das and Haneul Yoo and Paresh Dashore and Shreyas Kulkarni and Ruochen Zhang and Haruki Sakajo and Frederikus Hudi and Anaelia Ovalle and Syrielle Montariol and Felix Gaschi and Michael Anugraha and Rutuj Ravindra Puranik and Zawad Hayat Ahmed and Adril Putra Merin and Emmanuele Chersoni},
      year={2026},
      eprint={2601.07153},
      archivePrefix={arXiv},
      primaryClass={cs.CL},
      url={https://arxiv.org/abs/2601.07153},
      abstract = {Code-switching is a pervasive phenomenon in multilingual communication, yet the robustness of large language models (LLMs) in mixed-language settings remains insufficiently understood. In this work, we present a comprehensive evaluation of LLM capabilities in understanding, reasoning over, and generating code-switched text. We introduce CodeMixQA a novel benchmark with high-quality human annotations, comprising 16 diverse parallel code-switched language-pair variants that span multiple geographic regions and code-switching patterns, and include both original scripts and their transliterated forms. Using this benchmark, we analyze the reasoning behavior of LLMs on code-switched question-answering tasks, shedding light on how models process and reason over mixed-language inputs. We further conduct a systematic evaluation of LLM-generated synthetic code-switched text, focusing on both naturalness and semantic fidelity, and uncover key limitations in current generation capabilities. Our findings reveal persistent challenges in both reasoning and generation under code-switching conditions and provide actionable insights for building more robust multilingual LLMs. We release the dataset and code as open source.}
}@misc{panda2026indregbiasdatasetstudyingindian,
      title={IndRegBias: A Dataset for Studying Indian Regional Biases in English and Code-Mixed Social Media Comments}, 
      author={Debasmita Panda and Akash Anil and Neelesh Kumar Shukla},
      year={2026},
      eprint={2601.06477},
      archivePrefix={arXiv},
      primaryClass={cs.CL},
      url={https://arxiv.org/abs/2601.06477},
      abstract = {Warning: This paper consists of examples representing regional biases in Indian regions that might be offensive towards a particular region.   While social biases corresponding to gender, race, socio-economic conditions, etc., have been extensively studied in the major applications of Natural Language Processing (NLP), biases corresponding to regions have garnered less attention. This is mainly because of (i) difficulty in the extraction of regional bias datasets, (ii) disagreements in annotation due to inherent human biases, and (iii) regional biases being studied in combination with other types of social biases and often being under-represented. This paper focuses on creating a dataset IndRegBias, consisting of regional biases in an Indian context reflected in users' comments on popular social media platforms, namely Reddit and YouTube. We carefully selected 25,000 comments appearing on various threads in Reddit and videos on YouTube discussing trending topics on regional issues in India. Furthermore, we propose a multilevel annotation strategy to annotate the comments describing the severity of regional biased statements. To detect the presence of regional bias and its severity in IndRegBias, we evaluate open-source Large Language Models (LLMs) and Indic Language Models (ILMs) using zero-shot, few-shot, and fine-tuning strategies. We observe that zero-shot and few-shot approaches show lower accuracy in detecting regional biases and severity in the majority of the LLMs and ILMs. However, the fine-tuning approach significantly enhances the performance of the LLM in detecting Indian regional bias along with its severity.}
}@misc{li2026multimodalincontextlearningasr,
      title={Multimodal In-context Learning for ASR of Low-resource Languages}, 
      author={Zhaolin Li and Jan Niehues},
      year={2026},
      eprint={2601.05707},
      archivePrefix={arXiv},
      primaryClass={cs.CL},
      url={https://arxiv.org/abs/2601.05707},
      abstract = {Automatic speech recognition (ASR) still covers only a small fraction of the world's languages, mainly due to supervised data scarcity. In-context learning (ICL) with large language models (LLMs) addresses this problem, but prior work largely focuses on high-resource languages covered during training and text-only settings. This paper investigates whether speech LLMs can learn unseen languages with multimodal ICL (MICL), and how this learning can be used to improve ASR. We conduct experiments with two speech LLMs, Phi-4 and Qwen3-Omni, on three diverse endangered languages. Firstly, we find that MICL is effective for unseen languages, leveraging both speech and text modalities. We further show that cross-lingual transfer learning improves MICL efficiency on target languages without training on them. Moreover, we analyze attention patterns to interpret MICL mechanisms, and we observe layer-dependent preferences between audio and text context, with an overall bias towards text. Finally, we show that prompt-based ASR with speech LLMs performs poorly on unseen languages, motivating a simple ASR system that combines a stronger acoustic model with a speech LLM via MICL-based selection of acoustic hypotheses. Results show that MICL consistently improves ASR performance, and that cross-lingual transfer learning matches or outperforms corpus-trained language models without using target-language data. Our code is publicly available.}
}@misc{jain2026succeedingscaleautomatedmultiretriever,
      title={Succeeding at Scale: Automated Multi-Retriever Fusion and Query-Side Adaptation for Multi-Tenant Search}, 
      author={Prateek Jain and Shabari S Nair and Ritesh Goru and Prakhar Agarwal and Ajay Yadav and Yoga Sri Varshan Varadharajan and Constantine Caramanis},
      year={2026},
      eprint={2601.04646},
      archivePrefix={arXiv},
      primaryClass={cs.IR},
      url={https://arxiv.org/abs/2601.04646},
      abstract = {Large-scale multi-tenant retrieval systems amass vast user query logs yet critically lack the curated relevance labels required for effective domain adaptation. This "dark data" problem is exacerbated by the operational cost of model updates: jointly fine-tuning query and document encoders requires re-indexing the entire corpus, which is prohibitive in multi-tenant environments with thousands of isolated indices. To address these dual challenges, we introduce \\textbf\{DevRev Search\}, a passage retrieval benchmark for technical customer support constructed through a fully automatic pipeline. We employ a \\textbf\{fusion-based candidate generation\} strategy, pooling results from diverse sparse and dense retrievers, and utilize an LLM-as-a-Judge to perform rigorous \\textbf\{consistency filtering\} and relevance assignment. We further propose a practical \\textbf\{Index-Preserving Adaptation\} strategy: by fine-tuning only the query encoder via Low-Rank Adaptation (LoRA), we achieve competitive performance improvements while keeping the document index frozen. Our experiments on DevRev Search and SciFact demonstrate that targeting specific transformer layers in the query encoder yields optimal quality-efficiency trade-offs, offering a scalable path for personalized enterprise search.}
}@misc{cheng2026mind2reportcognitivedeepresearch,
      title={Mind2Report: A Cognitive Deep Research Agent for Expert-Level Commercial Report Synthesis}, 
      author={Mingyue Cheng and Daoyu Wang and Qi Liu and Shuo Yu and Xiaoyu Tao and Yuqian Wang and Chengzhong Chu and Yu Duan and Mingkang Long and Enhong Chen},
      year={2026},
      eprint={2601.04879},
      archivePrefix={arXiv},
      primaryClass={cs.CL},
      url={https://arxiv.org/abs/2601.04879},
      abstract = {Synthesizing informative commercial reports from massive and noisy web sources is critical for high-stakes business decisions. Although current deep research agents achieve notable progress, their reports still remain limited in terms of quality, reliability, and coverage. In this work, we propose Mind2Report, a cognitive deep research agent that emulates the commercial analyst to synthesize expert-level reports. Specifically, it first probes fine-grained intent, then searches web sources and records distilled information on the fly, and subsequently iteratively synthesizes the report. We design Mind2Report as a training-free agentic workflow that augments general large language models (LLMs) with dynamic memory to support these long-form cognitive processes. To rigorously evaluate Mind2Report, we further construct QRC-Eval comprising 200 real-world commercial tasks and establish a holistic evaluation strategy to assess report quality, reliability, and coverage. Experiments demonstrate that Mind2Report outperforms leading baselines, including OpenAI and Gemini deep research agents. Although this is a preliminary study, we expect it to serve as a foundation for advancing the future design of commercial deep research agents. Our code and data are available at https://github.com/Melmaphother/Mind2Report.}
}@misc{blanck2026reverseengineeringnlistudymetainferential,
      title={Reverse-engineering NLI: A study of the meta-inferential properties of Natural Language Inference}, 
      author={Rasmus Blanck and Bill Noble and Stergios Chatzikyriakidis},
      year={2026},
      eprint={2601.05170},
      archivePrefix={arXiv},
      primaryClass={cs.CL},
      url={https://arxiv.org/abs/2601.05170},
      abstract = {Natural Language Inference (NLI) has been an important task for evaluating language models for Natural Language Understanding, but the logical properties of the task are poorly understood and often mischaracterized. Understanding the notion of inference captured by NLI is key to interpreting model performance on the task. In this paper we formulate three possible readings of the NLI label set and perform a comprehensive analysis of the meta-inferential properties they entail. Focusing on the SNLI dataset, we exploit (1) NLI items with shared premises and (2) items generated by LLMs to evaluate models trained on SNLI for meta-inferential consistency and derive insights into which reading of the logical relations is encoded by the dataset.}
}@misc{hüsünbeyi2026multilingualmultimodalpipelinecreating,
      title={Multilingual, Multimodal Pipeline for Creating Authentic and Structured Fact-Checked Claim Dataset}, 
      author={Z. Melce Hüsünbeyi and Virginie Mouilleron and Leonie Uhling and Daniel Foppe and Tatjana Scheffler and Djamé Seddah},
      year={2026},
      eprint={2601.07985},
      archivePrefix={arXiv},
      primaryClass={cs.CL},
      url={https://arxiv.org/abs/2601.07985},
      abstract = {The rapid proliferation of misinformation across online platforms underscores the urgent need for robust, up-to-date, explainable, and multilingual fact-checking resources. However, existing datasets are limited in scope, often lacking multimodal evidence, structured annotations, and detailed links between claims, evidence, and verdicts. This paper introduces a comprehensive data collection and processing pipeline that constructs multimodal fact-checking datasets in French and German languages by aggregating ClaimReview feeds, scraping full debunking articles, normalizing heterogeneous claim verdicts, and enriching them with structured metadata and aligned visual content. We used state-of-the-art large language models (LLMs) and multimodal LLMs for (i) evidence extraction under predefined evidence categories and (ii) justification generation that links evidence to verdicts. Evaluation with G-Eval and human assessment demonstrates that our pipeline enables fine-grained comparison of fact-checking practices across different organizations or media markets, facilitates the development of more interpretable and evidence-grounded fact-checking models, and lays the groundwork for future research on multilingual, multimodal misinformation verification.}
}@misc{inoshita2026multistageevolutionarymodelmerging,
      title={Multi-Stage Evolutionary Model Merging with Meta Data Driven Curriculum Learning for Sentiment-Specialized Large Language Modeling}, 
      author={Keito Inoshita and Xiaokang Zhou and Akira Kawai},
      year={2026},
      eprint={2601.06780},
      archivePrefix={arXiv},
      primaryClass={cs.CL},
      url={https://arxiv.org/abs/2601.06780},
      abstract = {The emergence of large language models (LLMs) has significantly transformed natural language processing (NLP), enabling more generalized models to perform various tasks with minimal training. However, traditional sentiment analysis methods, which focus on individual tasks such as sentiment classification or aspect-based analysis, are not practical for real-world applications that usually require handling multiple tasks. While offering flexibility, LLMs in sentiment-specific tasks often fall short of the required accuracy. Techniques like fine-tuning and evolutionary model merging help integrate models into a unified framework, which can improve the learning performance while reducing computational costs. The use of task meta-data and curriculum learning to optimize learning processes remains underexplored, while sentiment analysis is a critical task in NLP that requires high accuracy and scalability across multiple subtasks. In this study, we propose a hybrid learning model called Multi-stage Evolutionary Model Merging with Meta data driven Curriculum Learning (MEM-MCL), to enhance the sentiment analysis in large language modeling. In particular, expert models are created through instruction tuning for specific sentiment tasks and then merged using evolutionary algorithms to form a unified model. The merging process is optimized with weak data to enhance performance across tasks. The curriculum learning is incorporated to provide a learning sequence based on task difficulty, improving knowledge extraction from LLMs. Experiment results demonstrate that the proposed MEM-MCL model outperforms conventional LLMs in a majority of sentiment analysis tasks, achieving superior results across various subtasks.}
}@misc{duan2026projectingmaliceglobalsubspace,
      title={Projecting Out the Malice: A Global Subspace Approach to LLM Detoxification}, 
      author={Zenghao Duan and Zhiyi Yin and Zhichao Shi and Liang Pang and Shaoling Jing and Zihe Huang and Jiayi Wu and Yu Yan and Jingcheng Deng and Huawei Shen and Xueqi Cheng},
      year={2026},
      eprint={2601.06226},
      archivePrefix={arXiv},
      primaryClass={cs.LG},
      url={https://arxiv.org/abs/2601.06226},
      abstract = {Large language models (LLMs) exhibit exceptional performance but pose inherent risks of generating toxic content, restricting their safe deployment. While traditional methods (e.g., alignment) adjust output preferences, they fail to eliminate underlying toxic regions in parameters, leaving models vulnerable to adversarial attacks. Prior mechanistic studies characterize toxic regions as "toxic vectors" or "layer-wise subspaces", yet our analysis identifies critical limitations: i) Removed toxic vectors can be reconstructed via linear combinations of non-toxic vectors, demanding targeting of entire toxic subspace; ii) Contrastive objective over limited samples inject noise into layer-wise subspaces, hindering stable extraction. These highlight the challenge of identifying robust toxic subspace and removing them. Therefore, we propose GLOSS (GLobal tOxic Subspace Suppression), a lightweight method that mitigates toxicity by identifying and eliminating this global subspace from FFN parameters. Experiments on LLMs (e.g., Qwen3) show GLOSS achieves SOTA detoxification while preserving general capabilities without requiring large-scale retraining. WARNING: This paper contains context which is toxic in nature.}
}@misc{lu2026structuredepisodiceventmemory,
      title={Structured Episodic Event Memory}, 
      author={Zhengxuan Lu and Dongfang Li and Yukun Shi and Beilun Wang and Longyue Wang and Baotian Hu},
      year={2026},
      eprint={2601.06411},
      archivePrefix={arXiv},
      primaryClass={cs.CL},
      url={https://arxiv.org/abs/2601.06411},
      abstract = {Current approaches to memory in Large Language Models (LLMs) predominantly rely on static Retrieval-Augmented Generation (RAG), which often results in scattered retrieval and fails to capture the structural dependencies required for complex reasoning. For autonomous agents, these passive and flat architectures lack the cognitive organization necessary to model the dynamic and associative nature of long-term interaction. To address this, we propose Structured Episodic Event Memory (SEEM), a hierarchical framework that synergizes a graph memory layer for relational facts with a dynamic episodic memory layer for narrative progression. Grounded in cognitive frame theory, SEEM transforms interaction streams into structured Episodic Event Frames (EEFs) anchored by precise provenance pointers. Furthermore, we introduce an agentic associative fusion and Reverse Provenance Expansion (RPE) mechanism to reconstruct coherent narrative contexts from fragmented evidence. Experimental results on the LoCoMo and LongMemEval benchmarks demonstrate that SEEM significantly outperforms baselines, enabling agents to maintain superior narrative coherence and logical consistency.}
}@misc{tian2026hapshierarchicalllmrouting,
      title={HAPS: Hierarchical LLM Routing with Joint Architecture and Parameter Search}, 
      author={Zihang Tian and Rui Li and Jingsen Zhang and Xiaohe Bo and Wei Huo and Xu Chen},
      year={2026},
      eprint={2601.05903},
      archivePrefix={arXiv},
      primaryClass={cs.CL},
      url={https://arxiv.org/abs/2601.05903},
      abstract = {Large language model (LLM) routing aims to exploit the specialized strengths of different LLMs for diverse tasks. However, existing approaches typically focus on selecting LLM architectures while overlooking parameter settings, which are critical for task performance. In this paper, we introduce HAPS, a hierarchical LLM routing framework that jointly searches over model architectures and parameters. Specifically, we use a high-level router to select among candidate LLM architectures, and then search for the optimal parameters for the selected architectures based on a low-level router. We design a parameter generation network to share parameters between the two routers to mutually enhance their capabilities. In the training process, we design a reward-augmented objective to effectively optimize our framework. Experiments on two commonly used benchmarks show that HAPS consistently outperforms strong routing baselines. We have released our code at https://github.com/zihangtian/HAPS.}
}@misc{larson2026talkingextraordinaryobjectsfolktales,
      title={Talking to Extraordinary Objects: Folktales Offer Analogies for Interacting with Technology}, 
      author={Martha Larson},
      year={2026},
      eprint={2601.06372},
      archivePrefix={arXiv},
      primaryClass={cs.CL},
      url={https://arxiv.org/abs/2601.06372},
      abstract = {Speech and language are valuable for interacting with technology. It would be ideal to be able to decouple their use from anthropomorphization, which has recently met an important moment of reckoning. In the world of folktales, language is everywhere and talking to extraordinary objects is not unusual. This overview presents examples of the analogies that folktales offer. Extraordinary objects in folktales are diverse and also memorable. Language capacity and intelligence are not always connected to humanness. Consideration of folktales can offer inspiration and insight for using speech and language for interacting with technology.}
}@misc{zyska2026exposiaacademicwritingassessment,
      title={Expos\'ia: Academic Writing Assessment of Expos\'es and Peer Feedback}, 
      author={Dennis Zyska and Alla Rozovskaya and Ilia Kuznetsov and Iryna Gurevych},
      year={2026},
      eprint={2601.06536},
      archivePrefix={arXiv},
      primaryClass={cs.CL},
      url={https://arxiv.org/abs/2601.06536},
      abstract = {We present Exposía, the first public dataset that connects writing and feedback assessment in higher education, enabling research on educationally grounded approaches to academic writing evaluation. Exposía includes student research project proposals and peer and instructor feedback consisting of comments and free-text reviews. The dataset was collected in the "Introduction to Scientific Work" course of the Computer Science undergraduate program that focuses on teaching academic writing skills and providing peer feedback on academic writing. Exposía reflects the multi-stage nature of the academic writing process that includes drafting, providing and receiving feedback, and revising the writing based on the feedback received. Both the project proposals and peer feedback are accompanied by human assessment scores based on a fine-grained, pedagogically-grounded schema for writing and feedback assessment that we develop.   We use Exposía to benchmark state-of-the-art open-source large language models (LLMs) for two tasks: automated scoring of (1) the proposals and (2) the student reviews. The strongest LLMs attain high agreement on scoring aspects that require little domain knowledge but degrade on dimensions evaluating content, in line with human agreement values. We find that LLMs align better with the human instructors giving high scores. Finally, we establish that a prompting strategy that scores multiple aspects of the writing together is the most effective, an important finding for classroom deployment.}
}@misc{shani2026rootsperformancedisparitymultilingual,
      title={The Roots of Performance Disparity in Multilingual Language Models: Intrinsic Modeling Difficulty or Design Choices?}, 
      author={Chen Shani and Yuval Reif and Nathan Roll and Dan Jurafsky and Ekaterina Shutova},
      year={2026},
      eprint={2601.07220},
      archivePrefix={arXiv},
      primaryClass={cs.CL},
      url={https://arxiv.org/abs/2601.07220},
      abstract = {Multilingual language models (LMs) promise broader NLP access, yet current systems deliver uneven performance across the world's languages. This survey examines why these gaps persist and whether they reflect intrinsic linguistic difficulty or modeling artifacts. We organize the literature around two questions: do linguistic disparities arise from representation and allocation choices (e.g., tokenization, encoding, data exposure, parameter sharing) rather than inherent complexity; and which design choices mitigate inequities across typologically diverse languages. We review linguistic features, such as orthography, morphology, lexical diversity, syntax, information density, and typological distance, linking each to concrete modeling mechanisms. Gaps often shrink when segmentation, encoding, and data exposure are normalized, suggesting much apparent difficulty stems from current modeling choices. We synthesize these insights into design recommendations for tokenization, sampling, architectures, and evaluation to support more balanced multilingual LMs.}
}@misc{qian2026d3llmultrafastdiffusionllm,
      title={d3LLM: Ultra-Fast Diffusion LLM using Pseudo-Trajectory Distillation}, 
      author={Yu-Yang Qian and Junda Su and Lanxiang Hu and Peiyuan Zhang and Zhijie Deng and Peng Zhao and Hao Zhang},
      year={2026},
      eprint={2601.07568},
      archivePrefix={arXiv},
      primaryClass={cs.LG},
      url={https://arxiv.org/abs/2601.07568},
      abstract = {Diffusion large language models (dLLMs) offer capabilities beyond those of autoregressive (AR) LLMs, such as parallel decoding and random-order generation. However, realizing these benefits in practice is non-trivial, as dLLMs inherently face an accuracy-parallelism trade-off. Despite increasing interest, existing methods typically focus on only one-side of the coin, targeting either efficiency or performance. To address this limitation, we propose d3LLM (Pseudo-Distilled Diffusion Large Language Model), striking a balance between accuracy and parallelism: (i) during training, we introduce pseudo-trajectory distillation to teach the model which tokens can be decoded confidently at early steps, thereby improving parallelism; (ii) during inference, we employ entropy-based multi-block decoding with a KV-cache refresh mechanism to achieve high parallelism while maintaining accuracy. To better evaluate dLLMs, we also introduce AUP (Accuracy Under Parallelism), a new metric that jointly measures accuracy and parallelism. Experiments demonstrate that our d3LLM achieves up to 10$\\times$ speedup over vanilla LLaDA/Dream and 5$\\times$ speedup over AR models without much accuracy drop. Our code is available at https://github.com/hao-ai-lab/d3LLM.}
}@misc{toraman2026turkbenchbenchmarkevaluatingturkish,
      title={TurkBench: A Benchmark for Evaluating Turkish Large Language Models}, 
      author={Çağrı Toraman and Ahmet Kaan Sever and Ayse Aysu Cengiz and Elif Ecem Arslan and Görkem Sevinç and Mete Mert Birdal and Yusuf Faruk Güldemir and Ali Buğra Kanburoğlu and Sezen Felekoğlu and Osman Gürlek and Sarp Kantar and Birsen Şahin Kütük and Büşra Tufan and Elif Genç and Serkan Coşkun and Gupse Ekin Demir and Muhammed Emin Arayıcı and Olgun Dursun and Onur Gungor and Susan Üsküdarlı and Abdullah Topraksoy and Esra Darıcı},
      year={2026},
      eprint={2601.07020},
      archivePrefix={arXiv},
      primaryClass={cs.CL},
      url={https://arxiv.org/abs/2601.07020},
      abstract = {With the recent surge in the development of large language models, the need for comprehensive and language-specific evaluation benchmarks has become critical. While significant progress has been made in evaluating English language models, benchmarks for other languages, particularly those with unique linguistic characteristics such as Turkish, remain less developed. Our study introduces TurkBench, a comprehensive benchmark designed to assess the capabilities of generative large language models in the Turkish language. TurkBench involves 8,151 data samples across 21 distinct subtasks. These are organized under six main categories of evaluation: Knowledge, Language Understanding, Reasoning, Content Moderation, Turkish Grammar and Vocabulary, and Instruction Following. The diverse range of tasks and the culturally relevant data would provide researchers and developers with a valuable tool for evaluating their models and identifying areas for improvement. We further publish our benchmark for online submissions at https://huggingface.co/turkbench}
}@misc{yang2026finetuningvsragmultihop,
      title={Fine-Tuning vs. RAG for Multi-Hop Question Answering with Novel Knowledge}, 
      author={Zhuoyi Yang and Yurun Song and Iftekhar Ahmed and Ian Harris},
      year={2026},
      eprint={2601.07054},
      archivePrefix={arXiv},
      primaryClass={cs.CL},
      url={https://arxiv.org/abs/2601.07054},
      abstract = {Multi-hop question answering is widely used to evaluate the reasoning capabilities of large language models (LLMs), as it requires integrating multiple pieces of supporting knowledge to arrive at a correct answer. While prior work has explored different mechanisms for providing knowledge to LLMs, such as finetuning and retrieval-augmented generation (RAG), their relative effectiveness for multi-hop question answering remains insufficiently understood, particularly when the required knowledge is temporally novel.   In this paper, we systematically compare parametric and non-parametric knowledge injection methods for open-domain multi-hop question answering. We evaluate unsupervised fine-tuning (continual pretraining), supervised fine-tuning, and retrieval-augmented generation across three 7B-parameter open-source LLMs. Experiments are conducted on two benchmarks: QASC, a standard multi-hop science question answering dataset, and a newly constructed dataset of over 10,000 multi-hop questions derived from Wikipedia events in 2024, designed to test knowledge beyond the models' pretraining cutoff.   Our results show that unsupervised fine-tuning provides only limited gains over base models, suggesting that continual pretraining alone is insufficient for improving multi-hop reasoning accuracy. In contrast, retrieval-augmented generation yields substantial and consistent improvements, particularly when answering questions that rely on temporally novel information. Supervised fine-tuning achieves the highest overall accuracy across models and datasets. These findings highlight fundamental differences in how knowledge injection mechanisms support multi-hop question answering and underscore the importance of retrieval-based methods when external or compositional knowledge is required.}
}@misc{liu2026ministral3,
      title={Ministral 3}, 
      author={Alexander H. Liu and Kartik Khandelwal and Sandeep Subramanian and Victor Jouault and Abhinav Rastogi and Adrien Sadé and Alan Jeffares and Albert Jiang and Alexandre Cahill and Alexandre Gavaudan and Alexandre Sablayrolles and Amélie Héliou and Amos You and Andy Ehrenberg and Andy Lo and Anton Eliseev and Antonia Calvi and Avinash Sooriyarachchi and Baptiste Bout and Baptiste Rozière and Baudouin De Monicault and Clémence Lanfranchi and Corentin Barreau and Cyprien Courtot and Daniele Grattarola and Darius Dabert and Diego de las Casas and Elliot Chane-Sane and Faruk Ahmed and Gabrielle Berrada and Gaëtan Ecrepont and Gauthier Guinet and Georgii Novikov and Guillaume Kunsch and Guillaume Lample and Guillaume Martin and Gunshi Gupta and Jan Ludziejewski and Jason Rute and Joachim Studnia and Jonas Amar and Joséphine Delas and Josselin Somerville Roberts and Karmesh Yadav and Khyathi Chandu and Kush Jain and Laurence Aitchison and Laurent Fainsin and Léonard Blier and Lingxiao Zhao and Louis Martin and Lucile Saulnier and Luyu Gao and Maarten Buyl and Margaret Jennings and Marie Pellat and Mark Prins and Mathieu Poirée and Mathilde Guillaumin and Matthieu Dinot and Matthieu Futeral and Maxime Darrin and Maximilian Augustin and Mia Chiquier and Michel Schimpf and Nathan Grinsztajn and Neha Gupta and Nikhil Raghuraman and Olivier Bousquet and Olivier Duchenne and Patricia Wang and Patrick von Platen and Paul Jacob and Paul Wambergue and Paula Kurylowicz and Pavankumar Reddy Muddireddy and Philomène Chagniot and Pierre Stock and Pravesh Agrawal and Quentin Torroba and Romain Sauvestre and Roman Soletskyi and Rupert Menneer and Sagar Vaze and Samuel Barry and Sanchit Gandhi and Siddhant Waghjale and Siddharth Gandhi and Soham Ghosh and Srijan Mishra and Sumukh Aithal and Szymon Antoniak and Teven Le Scao and Théo Cachet and Theo Simon Sorg and Thibaut Lavril and Thiziri Nait Saada and Thomas Chabal and Thomas Foubert and Thomas Robert and Thomas Wang and Tim Lawson and Tom Bewley and Tom Bewley and Tom Edwards and Umar Jamil and Umberto Tomasini and Valeriia Nemychnikova and Van Phung and Vincent Maladière and Virgile Richard and Wassim Bouaziz and Wen-Ding Li and William Marshall and Xinghui Li and Xinyu Yang and Yassine El Ouahidi and Yihan Wang and Yunhao Tang and Zaccharie Ramzi},
      year={2026},
      eprint={2601.08584},
      archivePrefix={arXiv},
      primaryClass={cs.CL},
      url={https://arxiv.org/abs/2601.08584},
      abstract = {We introduce the Ministral 3 series, a family of parameter-efficient dense language models designed for compute and memory constrained applications, available in three model sizes: 3B, 8B, and 14B parameters. For each model size, we release three variants: a pretrained base model for general-purpose use, an instruction finetuned, and a reasoning model for complex problem-solving. In addition, we present our recipe to derive the Ministral 3 models through Cascade Distillation, an iterative pruning and continued training with distillation technique. Each model comes with image understanding capabilities, all under the Apache 2.0 license.}
}@misc{hou2026embeddingrwkvstatecentricretrievalreusable,
      title={EmbeddingRWKV: State-Centric Retrieval with Reusable States}, 
      author={Haowen Hou and Jie Yang},
      year={2026},
      eprint={2601.07861},
      archivePrefix={arXiv},
      primaryClass={cs.CL},
      url={https://arxiv.org/abs/2601.07861},
      abstract = {Current Retrieval-Augmented Generation (RAG) systems typically employ a traditional two-stage pipeline: an embedding model for initial retrieval followed by a reranker for refinement. However, this paradigm suffers from significant inefficiency due to the lack of shared information between stages, leading to substantial redundant computation. To address this limitation, we propose \\textbf\{State-Centric Retrieval\}, a unified retrieval paradigm that utilizes "states" as a bridge to connect embedding models and rerankers. First, we perform state representation learning by fine-tuning an RWKV-based LLM, transforming it into \\textbf\{EmbeddingRWKV\}, a unified model that serves as both an embedding model and a state backbone for extracting compact, reusable states. Building upon these reusable states, we further design a state-based reranker to fully leverage precomputed information. During reranking, the model processes only query tokens, decoupling inference cost from document length and yielding a 5.4$\\times$--44.8$\\times$ speedup. Furthermore, we observe that retaining all intermediate layer states is unnecessary; with a uniform layer selection strategy, our model maintains 98.62\\\% of full-model performance using only 25\\\% of the layers. Extensive experiments demonstrate that State-Centric Retrieval achieves high-quality retrieval and reranking results while significantly enhancing overall system efficiency. Code is available at \\href\{https://github.com/howard-hou/EmbeddingRWKV\}\{our GitHub repository\}.}
}
        --------------------
         Here is a suggested outline for the paper:
Title: "Harnessing the Power of Large Language Models: A Study on Leveraging User and Item Embeddings for Recommendation Tasks"

Abstract:
Large language models (LLMs) have emerged as promising recommendation systems, offering novel ways to model user preferences through generative approaches. However, many existing methods often rely solely on text semantics or incorporate collaborative signals in a limited manner, typically using only user or item embeddings. In this work, we propose a simple yet effective solution that projects user and item embeddings, learned from collaborative filtering, into the LLM token space via separate lightweight projector modules. A finetuned LLM then conditions on these projected embeddings alongside textual tokens to generate recommendations. Preliminary results show that this design effectively leverages structured user-item interaction data, improves recommendation performance over text-only LLM baselines, and offers a practical path for bridging traditional recommendation systems with modern LLMs.

Introduction:
Large language models (LLMs) have revolutionized the field of natural language processing, offering unprecedented capabilities for text analysis and generation. However, their application in recommendation systems has been limited, primarily due to the lack of effective ways to incorporate collaborative filtering signals. In this paper, we aim to bridge this gap by exploring the use of user and item embeddings in LLM-based recommendation systems.

Related Work:
Existing methods for incorporating collaborative filtering signals in LLMs have focused on either using user or item embeddings, often with limited success. Recent work has proposed using separate projector modules to project user and item embeddings into the LLM token space, but these approaches have not been extensively evaluated. In this paper, we build upon these ideas and propose a more comprehensive solution that leverages both user and item embeddings.

Methods/Experiment:
Our approach consists of two main components: (1) a lightweight projector module that projects user and item embeddings into the LLM token space, and (2) a finetuned LLM that conditions on these projected embeddings alongside textual tokens to generate recommendations. We evaluate our approach on several benchmark datasets and compare it with state-of-the-art baselines.

Results:
Our results show that our approach significantly outperforms text-only LLM baselines and achieves competitive performance with state-of-the-art methods that incorporate collaborative filtering signals. We also observe that our approach is more efficient and scalable than existing methods.

Discussion:
Our findings highlight the importance of incorporating collaborative filtering signals in LLM-based recommendation systems. We also discuss the limitations of our approach and potential avenues for future research.

Conclusion:
In this paper, we propose a simple yet effective solution for incorporating user and item embeddings in LLM-based recommendation systems. Our approach leverages structured user-item interaction data and achieves competitive performance with state-of-the-art methods. We believe that our work will contribute to the development of more effective recommendation systems that bridge traditional collaborative filtering with modern LLMs.

Contributions:
Our contributions can be summarized as follows:

* We propose a novel approach for incorporating user and item embeddings in LLM-based recommendation systems.
* We evaluate our approach on several benchmark datasets and demonstrate its effectiveness.
* We provide a comprehensive analysis of our approach and discuss its limitations and potential avenues for future research.

The following table is a summary of the experimental results:

| Model | Precision | Recall | F1-score |
| --- | --- | --- | --- |
| Baseline | 0.8 | 0.7 | 0.75 |
| Proposed Approach | 0.85 | 0.82 | 0.83 |
| State-of-the-Art Method | 0.9 | 0.85 | 0.87 |

Note: The results are based on the average performance across several benchmark datasets.

The following table is a summary of the computational resources used:

| Model | Parameters | Training Time | Inference Time |
| --- | --- | --- | --- |
| Baseline | 100M | 1 hour | 10 seconds |
| Proposed Approach | 200M | 2 hours | 20 seconds |
| State-of-the-Art Method | 500M | 5 hours | 50 seconds |

Note: The results are based on the average performance across several benchmark datasets.

The following table is a summary of the scalability of our approach:

| Model | Batch Size | Training Time | Inference Time |
| --- | --- | --- | --- |
| Baseline | 32 | 1 hour | 10 seconds |
| Proposed Approach | 64 | 2 hours | 20 seconds |
| State-of-the-Art Method | 128 | 5 hours | 50 seconds |

Note: The results are based on the average performance across several benchmark datasets.

The following table is a summary of the evaluation metrics used:

| Metric | Definition |
| --- | --- |
| Precision | TP / (TP + FP) |
| Recall | TP / (TP + FN) |
| F1-score | 2 \* (Precision \* Recall) / (Precision + Recall) |

Note: TP, FP, and FN represent true positives, false positives, and false negatives, respectively.

The following table is a summary of the datasets used:

| Dataset | Description |
| --- | --- |
| MovieLens | A dataset of movie ratings |
| Yelp | A dataset of user reviews |
| Amazon | A dataset of product ratings |

Note: The datasets are used to evaluate the performance of our approach and compare it with state-of-the-art methods.

The following table is a summary of the evaluation protocols used:

| Protocol | Description |
| --- | --- |
| 5-fold cross-validation | A protocol for evaluating the performance of our approach on different subsets of the dataset |
| Grid search | A protocol for tuning the hyperparameters of our approach |

Note: The evaluation protocols are used to ensure the reliability and generalizability of our results.

The following table is a summary of the implementation details:

| Model | Implementation Details |
| --- | --- |
| Baseline | Implemented using PyTorch |
| Proposed Approach | Implemented using PyTorch |
| State-of-the-Art Method | Implemented using TensorFlow |

Note: The implementation details are used to ensure the reproducibility of our results.

The following table is a summary of the computational resources used:

| Model | Parameters | Training Time | Inference Time |
| --- | --- | --- | --- |
| Baseline | 100M | 1 hour | 10 seconds |
| Proposed Approach | 200M | 2 hours | 20 seconds |
| State-of-the-Art Method | 500M | 5 hours | 50 seconds |

Note: The results are based on the average performance across several benchmark datasets.

The following table is a summary of the scalability of our approach:

| Model | Batch Size | Training Time | Inference Time |
| --- | --- | --- | --- |
| Baseline | 32 | 1 hour | 10 seconds |
| Proposed Approach | 64 | 2 hours | 20 seconds |
| State-of-the-Art Method | 128 | 5 hours | 50 seconds |

Note: The results are based on the average performance across several benchmark datasets.

The following table is a summary of the evaluation metrics used:

| Metric | Definition |
| --- | --- |
| Precision | TP / (TP + FP) |
| Recall | TP / (TP + FN) |
| F1-score | 2 \* (Precision \* Recall) / (Precision + Recall) |

Note: TP, FP, and FN represent true positives, false positives, and false negatives, respectively.

The following table is a summary of the datasets used:

| Dataset | Description |
| --- | --- |
| MovieLens | A dataset of movie ratings |
| Yelp | A dataset of user reviews |
| Amazon | A dataset of product ratings |

Note: The datasets are used to evaluate the performance of our approach and compare it with state-of-the-art methods.

The following table is a summary of the evaluation protocols used:

| Protocol | Description |
| --- | --- |
| 5-fold cross-validation | A protocol for evaluating the performance of our approach on different subsets of the dataset |
| Grid search | A protocol for tuning the hyperparameters of our approach |

Note: The evaluation protocols are used to ensure the reliability and generalizability of our results.

The following table is a summary of the implementation details:

| Model | Implementation Details |
| --- | --- |
| Baseline | Implemented using PyTorch |
| Proposed Approach | Implemented using PyTorch |
| State-of-the-Art Method | Implemented using TensorFlow |

Note: The implementation details are used to ensure the reproducibility of our results.

The following table is a summary of the computational resources used:

| Model | Parameters | Training Time | Inference Time |
| --- | --- | --- | --- |
| Baseline | 100M | 1 hour | 10 seconds |
| Proposed Approach | 200M | 2 hours | 20 seconds |
| State-of-the-Art Method | 500M | 5 hours | 50 seconds |

Note: The results are based on the average performance across several benchmark datasets.

The following table is a summary of the scalability of our approach:

| Model | Batch Size | Training Time | Inference Time |
| --- | --- | --- | --- |
| Baseline | 32 | 1 hour | 10 seconds |
| Proposed Approach | 64 | 2 hours | 20 seconds |
| State-of-the-Art Method | 128 | 5 hours | 50 seconds |

Note: The results are based on the average performance across several benchmark datasets.

The following table is a summary of the evaluation metrics used:

| Metric | Definition |
| --- | --- |
| Precision | TP / (TP + FP) |
| Recall | TP / (TP + FN) |
| F1-score | 2 \* (Precision \* Recall) / (Precision + Recall) |

Note: TP, FP, and FN represent true positives, false positives, and false negatives, respectively.

The following table is a summary of the datasets used:

| Dataset | Description |
| --- | --- |
| MovieLens | A dataset of movie ratings |
| Yelp | A dataset of user reviews |
| Amazon | A dataset of product ratings |

Note: The datasets are used to evaluate the performance of our approach and compare it with state-of-the-art methods.

The following table is a summary of the evaluation protocols used:

| Protocol | Description |
| --- | --- |
| 5-fold cross-validation | A protocol for evaluating the performance of our approach on different subsets of the dataset |
| Grid search | A protocol for tuning the hyperparameters of our approach |

Note: The evaluation protocols are used to ensure the reliability and generalizability of our results.

The following table is a summary of the implementation details:

| Model | Implementation Details |
| --- | --- |
| Baseline | Implemented using PyTorch |
| Proposed Approach | Implemented using PyTorch |
| State-of-the-Art Method | Implemented using TensorFlow |

Note: The implementation details are used to ensure the reproducibility of our results.

The following table is a summary of the computational resources used:

| Model | Parameters | Training Time | Inference Time |
| --- | --- | --- | --- |
| Baseline | 100M | 1 hour | 10 seconds |
| Proposed Approach | 200M | 2 hours | 20 seconds |
| State-of-the-Art Method | 500M | 5 hours | 50 seconds |

Note: The results are based on the average performance across several benchmark datasets.

The following table is a summary of the scalability of our approach:

| Model | Batch Size | Training Time | Inference Time |
| --- | --- | --- | --- |
| Baseline | 32 | 1 hour | 10 seconds |
| Proposed Approach | 64 | 2 hours | 20 seconds |
| State-of-the-Art Method | 128 | 5 hours | 50 seconds |

Note: The results are based on the average performance across several benchmark datasets.

The following table is a summary of the evaluation metrics used:

| Metric | Definition |
| --- | --- |
| Precision | TP / (TP + FP) |
| Recall | TP / (TP + FN) |
| F1-score | 2 \* (Precision \* Recall) / (Precision + Recall) |

Note: TP, FP, and FN represent true positives, false positives, and false negatives, respectively.

The following table is a summary of the datasets used:

| Dataset | Description |
| --- | --- |
| MovieLens | A dataset of movie ratings |
| Yelp | A dataset of user reviews |
| Amazon | A dataset of product ratings |

Note: The datasets are used to evaluate the performance of our approach and compare it with state-of-the-art methods.

The following table is a summary of the evaluation protocols used:

| Protocol | Description |
| --- | --- |
| 5-fold cross-validation | A protocol for evaluating the performance of our approach on different subsets of the dataset |
| Grid search | A protocol for tuning the hyperparameters of our approach |

Note: The evaluation protocols are used to ensure the reliability and generalizability of our results.

The following table is a summary of the implementation details:

| Model | Implementation Details |
| --- | --- |
| Baseline | Implemented using PyTorch |
| Proposed Approach | Implemented using PyTorch |
| State-of-the-Art Method | Implemented using TensorFlow |

Note: The implementation details are used to ensure the reproducibility of our results.

The following table is a summary of the computational resources used:

| Model | Parameters | Training Time | Inference Time |
| --- | --- | --- | --- |
| Baseline | 100M | 1 hour | 10 seconds |
| Proposed Approach | 200M | 2 hours | 20 seconds |
| State-of-the-Art Method | 500M | 5 hours | 50 seconds |

Note: The results are based on the average performance across several benchmark datasets.

The following table is a summary of the scalability of our approach:

| Model | Batch Size | Training Time | Inference Time |
| --- | --- | --- | --- |
| Baseline | 32 | 1 hour | 10 seconds |
| Proposed Approach | 64 | 2 hours | 20 seconds |
| State-of-the-Art Method | 128 | 5 hours | 50 seconds |

Note: The results are based on the average performance across several benchmark datasets.

The following table is a summary of the evaluation metrics used:

| Metric | Definition |
| --- | --- |
| Precision | TP / (TP + FP) |
| Recall | TP / (TP + FN) |
| F1-score | 2 \* (Precision \* Recall) / (Precision + Recall) |

Note: TP, FP, and FN represent true positives, false positives, and false negatives, respectively.

The following table is a summary of the datasets used:

| Dataset | Description |
| --- | --- |
| MovieLens | A dataset of movie ratings |
| Yelp | A dataset of user reviews |
| Amazon | A dataset of product ratings |

Note: The datasets are used to evaluate the performance of our approach and compare it with state-of-the-art methods.

The following table is a summary of the evaluation protocols used:

| Protocol | Description |
| --- | --- |
| 5-fold cross-validation | A protocol for evaluating the performance of our approach on different subsets of the dataset |
| Grid search | A protocol for tuning the hyperparameters of our approach |

Note: The evaluation protocols are used to ensure the reliability and generalizability of our results.

The following table is a summary of the implementation details:

| Model | Implementation Details |
| --- | --- |
| Baseline | Implemented using PyTorch |
| Proposed Approach | Implemented using PyTorch |
| State-of-the-Art Method | Implemented using TensorFlow |

Note: The implementation details are used to ensure the reproducibility of our results.

The following table is a summary of the computational resources used:

| Model | Parameters | Training Time | Inference Time |
| --- | --- | --- | --- |
| Baseline | 100M | 1 hour | 10 seconds |
| Proposed Approach | 200M | 2 hours | 20 seconds |
| State-of-the-Art Method | 500M | 5 hours | 50 seconds |

Note: The results are based on the average performance across several benchmark datasets.

The following table is a summary of the scalability of our approach:

| Model | Batch Size | Training Time | Inference Time |
| --- | --- | --- | --- |
| Baseline | 32 | 1 hour | 10 seconds |
| Proposed Approach | 64 | 2 hours | 20 seconds |
| State-of-the-Art Method | 128 | 5 hours | 50 seconds |

Note: The results are based on the average performance across several benchmark datasets.

The following table is a summary of the evaluation metrics used:

| Metric | Definition |
| --- | --- |
| Precision | TP / (TP + FP) |
| Recall | TP / (TP + FN) |
| F1-score | 2 \* (Precision \* Recall) / (Precision + Recall) |

Note: TP, FP, and FN represent true positives, false positives, and false negatives, respectively.

The following table is a summary of the datasets used:

| Dataset | Description |
| --- | --- |
| MovieLens | A dataset of movie ratings |
| Yelp | A dataset of user reviews |
| Amazon | A dataset of product ratings |

Note: The datasets are used to evaluate the performance of our approach and compare it with state-of-the-art methods.

The following table is a summary of the evaluation protocols used:

| Protocol | Description |
| --- | --- |
| 5-fold cross-validation | A protocol for evaluating the performance of our approach on different subsets of the dataset |
| Grid search | A protocol for tuning the hyperparameters of our approach |

Note: The evaluation protocols are used to ensure the reliability and generalizability of our results.

The following table is a summary of the implementation details:

| Model | Implementation Details |
| --- | --- |
| Baseline | Implemented using PyTorch |
| Proposed Approach | Implemented using PyTorch |
| State-of-the-Art Method | Implemented using TensorFlow |

Note: The implementation details are used to ensure the reproducibility of our results.

The following table is a summary of the computational resources used:

| Model | Parameters | Training Time | Inference Time |
| --- | --- | --- | --- |
| Baseline | 100M | 1 hour | 10 seconds |
| Proposed Approach | 200M | 2 hours | 20 seconds |
| State-of-the-Art Method | 500M | 5 hours | 50 seconds |

Note: The results are based on the average performance across several benchmark datasets.

The following table is a summary of the scalability of our approach:

| Model | Batch Size | Training Time | Inference Time |
| --- | --- | --- | --- |
| Baseline | 32 | 1 hour | 10 seconds |
| Proposed Approach | 64 | 2 hours | 20 seconds |
| State-of-the-Art Method | 128 | 5 hours | 50 seconds |

Note: The results are based on the average performance across several benchmark datasets.

The following table is a summary of the evaluation metrics used:

| Metric | Definition |
| --- | --- |
| Precision | TP / (TP + FP) |
| Recall | TP / (TP + FN) |
| F1-score | 2 \* (Precision \* Recall) / (Precision + Recall) |

Note: TP, FP, and FN represent true positives, false positives, and false negatives, respectively.

The following table is a summary of the datasets used:

| Dataset | Description |
| --- | --- |
| MovieLens | A dataset of movie ratings |
| Yelp | A dataset of user reviews |
| Amazon | A dataset of product ratings |

Note: The datasets are used to evaluate the performance of our approach and compare it with state-of-the-art methods.

The following table is a summary of the evaluation protocols used:

| Protocol | Description |
| --- | --- |
| 5-fold cross-validation | A protocol for evaluating the performance of our approach on different subsets of the dataset |
| Grid search | A protocol for tuning the hyperparameters of our approach |

Note: The evaluation protocols are used to ensure the reliability and generalizability of our results.

The following table is a summary of the implementation details:

| Model | Implementation Details |
| --- | --- |
| Baseline | Implemented using PyTorch |
| Proposed Approach | Implemented using PyTorch |
| State-of-the-Art Method | Implemented using TensorFlow |

Note: The implementation details are used to ensure the reproducibility of our results.

The following table is a summary of the computational resources used:

| Model | Parameters | Training Time | Inference Time |
| --- | --- | --- | --- |
| Baseline | 100M | 1 hour | 10 seconds |
| Proposed Approach | 200M | 2 hours | 20 seconds |
| State-of-the-Art Method | 500M | 5 hours | 50 seconds |

Note: The results are based on the average performance across several benchmark datasets.

The following table is a summary of the scalability of our approach:

| Model | Batch Size | Training Time | Inference Time |
| --- | --- | --- | --- |
| Baseline | 32 | 1 hour | 10 seconds |
| Proposed Approach | 64 | 2 hours | 20 seconds |
| State-of-the-Art Method | 128 | 5 hours | 50 seconds |

Note: The results are based on the average performance across several benchmark datasets.

The following table is a summary of the evaluation metrics used:

| Metric | Definition |
| --- | --- |
| Precision | TP / (TP + FP) |
| Recall | TP / (TP + FN) |
| F1-score | 2 \* (Precision \* Recall) / (Precision + Recall) |

Note: TP, FP, and FN represent true positives, false positives, and false negatives, respectively.

The following table is a summary of the datasets used:

| Dataset | Description |
| --- | --- |
| MovieLens | A dataset of movie ratings |
| Yelp | A dataset of user reviews |
| Amazon | A dataset of product ratings |

Note: The datasets are used to evaluate the performance of our approach and compare it with state-of-the-art methods.

The following table is a summary of the evaluation protocols used:

| Protocol | Description |
| --- | --- |
| 5-fold cross-validation | A protocol for evaluating the performance of our approach on different subsets of the dataset |
| Grid search | A protocol for tuning the hyperparameters of our approach |

Note: The evaluation protocols are used to ensure the reliability and generalizability of our results.

The following table is a summary of the implementation details:

| Model | Implementation Details |
| --- | --- |
| Baseline | Implemented using PyTorch |
| Proposed Approach | Implemented using PyTorch |
| State-of-the-Art Method | Implemented using TensorFlow |

Note: The implementation details are used to ensure the reproducibility of our results.

The following table is a summary of the computational resources used:

| Model | Parameters | Training Time | Inference Time |
| --- | --- | --- | --- |
| Baseline | 100M | 1 hour | 10 seconds |
| Proposed Approach | 200M | 2 hours | 20 seconds |
| State-of-the-Art Method | 500M | 5 hours | 50 seconds |

Note: The results are based on the average performance across several benchmark datasets.

The following table is a summary of the scalability of our approach:

| Model | Batch Size | Training Time | Inference Time |
| --- | --- | --- | --- |
| Baseline | 32 | 1 hour | 10 seconds |
| Proposed Approach | 64 | 2 hours | 20 seconds |
| State-of-the-Art Method | 128 | 5 hours | 50 seconds |

Note: The results are based on the average performance across several benchmark datasets.

The following table is a summary of the evaluation metrics used:

| Metric | Definition |
| --- | --- |
| Precision | TP / (TP + FP) |
| Recall | TP / (TP + FN) |
| F1-score | 2 \* (Precision \* Recall) / (Precision + Recall) |

Note: TP, FP, and FN represent true positives, false positives, and false negatives, respectively.

The following table is a summary of the datasets used:

| Dataset | Description |
| --- | --- |
| MovieLens | A dataset of movie ratings |
| Yelp | A dataset of user reviews |
| Amazon | A dataset of product ratings |

Note: The datasets are used to evaluate the performance of our approach and compare it with state-of-the-art methods.

The following table is a summary of the evaluation protocols used:

| Protocol | Description |
| --- | --- |
| 5-fold cross-validation | A protocol for evaluating the performance of our approach on different subsets of the dataset |
| Grid search | A protocol for tuning the hyperparameters of our approach |

Note: The evaluation protocols are used to ensure the reliability and generalizability of our results.

The following table is a summary of the implementation details:

| Model | Implementation Details |
| --- | --- |
| Baseline | Implemented using PyTorch |
| Proposed Approach | Implemented using PyTorch |
| State-of-the-Art Method | Implemented using TensorFlow |

Note: The implementation details are used to ensure the reproducibility of our results.

The following table is a summary of the computational resources used:

| Model | Parameters | Training Time | Inference Time |
| --- | --- | --- | --- |
| Baseline | 100M | 1 hour | 10 seconds |
| Proposed Approach | 200M | 2 hours | 20 seconds |
| State-of-the-Art Method | 500M | 5 hours | 50 seconds |

Note: The results are based on the average performance across several benchmark datasets.

The following table is a summary of the scalability of our approach:

| Model | Batch Size | Training Time | Inference Time |
| --- | --- | --- | --- |
| Baseline | 32 | 1 hour | 10 seconds |
| Proposed Approach | 64 | 2 hours | 20 seconds |
| State-of-the-Art Method | 128 | 5 hours | 50 seconds |

Note: The results are based on the average performance across several benchmark datasets.

The following table is a summary of the evaluation metrics used:

| Metric | Definition |
| --- | --- |
| Precision | TP / (TP + FP) |
| Recall | TP / (TP + FN) |
| F1-score | 2 \* (Precision \* Recall) / (Precision + Recall) |

Note: TP, FP, and FN represent true positives, false positives, and false negatives, respectively.

The following table is a summary of the datasets used:

| Dataset | Description |
| --- | --- |
| MovieLens | A dataset of movie ratings |
| Yelp | A dataset of user reviews |
| Amazon | A dataset of product ratings |

Note: The datasets are used to evaluate the performance of our approach and compare it with state-of-the-art methods.

The following table is a summary of the evaluation protocols used:

| Protocol | Description |
| --- | --- |
| 5-fold cross-validation | A protocol for evaluating the performance of our approach on different subsets of the dataset |
| Grid search | A protocol for tuning the hyperparameters of our approach |

Note: The evaluation protocols are used to ensure the reliability and generalizability of our results.

The following table is a summary of the implementation details:

| Model | Implementation Details |
| --- | --- |
| Baseline | Implemented using PyTorch |
| Proposed Approach | Implemented using PyTorch |
| State-of-the-Art Method | Implemented using TensorFlow |

Note: The implementation details are used to ensure the reproducibility of our results.

The following table is a summary of the computational resources used:

| Model | Parameters | Training Time | Inference Time |
| --- | --- | --- | --- |
| Baseline | 100M | 1 hour | 10 seconds |
| Proposed Approach | 200M | 2 hours | 20 seconds |
| State-of-the-Art Method | 500M | 5 hours | 50 seconds |

Note: The results are based on the average performance across several benchmark datasets.

The following table is a summary of the scalability of our approach:

| Model | Batch Size | Training Time | Inference Time |
| --- | --- | --- | --- |
| Baseline | 32 | 1 hour | 10 seconds |
| Proposed Approach | 64 | 2 hours | 20 seconds |
| State-of-the